% sec-02: execution and settlement.  Table 24 and the actionbox.
\section{Execution and Settlement}

Six lines are entered in an order set by an external constraint rather than by
preference. The auction cutoff is a hard stop and a limit order is not, so the
auction bids go first.

\begin{apptable}
\begin{tabularx}{\textwidth}{>{\raggedright\arraybackslash}p{0.5cm} >{\raggedright\arraybackslash}p{3.4cm} >{\raggedright\arraybackslash}p{1.2cm} >{\raggedright\arraybackslash}p{2.8cm} >{\raggedright\arraybackslash}X}
\headrow \thead{\#} & \thead{Instrument} & \thead{Size} & \thead{Entry} & \thead{Settlement basis, and what is recorded}\\
\midrule
1 & Treasury bill, about 3 months & 20 pct & First, auction & Auction settlement date on the announcement; CUSIP, par, discount price, maturity \\
2 & Treasury bill, about 6 months & 20 pct & First, auction & The same \\
3 & Treasury bill, about 9 months & 20 pct & First, auction & The same \\
4 & Treasury note or bill, about 12 months & 20 pct & Second, limit, day & Regular way secondary; CUSIP, par, price, trade and settlement dates \\
5 & Short-duration Treasury ETF & 15 pct & Third, limit, day, mid-session & Standard fund settlement; ticker, shares, execution price, dates \\
6 & Government money market sweep & 5 pct & Automatic & Overnight; balance only \\
\end{tabularx}
\end{apptable}
\vspace{-0.60cm}
\tabcap{Table 24. Six lines with size, entry order, and settlement basis, in the\\
sequence an external auction cutoff imposes rather than one chosen freely.}

\subsection{The record written on the day of the fill}

Not at quarter end. A ledger written from statements three months later
reconciles to the statements and to nothing else. Every line records its
instrument identifier, side and quantity, price, trade date, settlement date,
maturity date where it has one, the share of reserve it represented on the
morning of entry, and a confirmation that auto-reinvestment is off.

Treasury bills are issued at a discount and the accretion is interest income for
federal purposes recognized at maturity or sale. Interest on direct United States
Treasury obligations is exempt from California personal income tax and the
exemption flows through a single-member limited liability company treated as a
disregarded entity, subject to confirmation by the company's tax preparer
\cite{treasuryyield2026}.

\subsection{If a line does not fill}

An auction line that misses a cutoff is carried to the next auction. It is not
replaced by a secondary purchase at a worse basis in order to feel finished. A
limit line that is not reached is canceled at the close and re-priced the next
morning, and \textbf{is never converted to a market order to complete a day}.

That last instruction is the most important line in this section. A ladder built
to avoid being forced should not be completed by forcing it, and the temptation
to do so arrives precisely on the day the rest of the block went well.

\subsection{Reconciliation at the close}

Every line either filled or is recorded as unfilled with a reason, so no line is
unaccounted for. The six shares sum to one hundred percent of the reserve as of
this morning. No position outside the authorized set exists in the account, which
is a control question rather than a preference question. Auto-reinvestment is off
on every rung, confirmed in the account settings rather than in an order ticket.
And there is no margin balance, no option position, and no lending consent.

\subsection{Why the entry order is not a preference}

An auction cutoff is external, published, and unforgiving: a non-competitive bid
submitted after a broker's internal cutoff is not submitted at all, and the next
opportunity is the following auction. A limit order is none of those things; it
rests through the session and can be canceled at the close and re-priced
tomorrow.

Sequencing therefore follows the constraint rather than the convenience. The
three auction lines go first because they are the only lines whose window can
close while the rest of the session is still open, and the exchange-traded line
goes last because its only timing constraint is internal, which is that it is not
entered in the first or last fifteen minutes of the session.

The general rule is worth naming because it recurs: where one action has an
external deadline and another does not, the external one goes first even when the
other is more urgent to the person doing it.

The rule has a cost, and the cost is that the most consequential action of a
given day is frequently not the first one taken. That inversion is uncomfortable
and it is correct. Urgency is a property of the person holding the list, while a
cutoff is a property of the world, and only one of the two moves when a founder
decides it should. Writing the sequence down the night before is what keeps the
distinction intact at nine in the morning, when the temptation is to begin with
whatever felt largest at breakfast.

\subsection{The single approval step}

\actionbox{One decision is open on this day}{Approve the release, confirming that
nothing changed overnight that would invalidate the preceding day's authorization,
and adopt the weekly cadence in \S4. The release is mechanical and is governed by
the six checks in \S1; the cadence is the decision that makes the next block a
substitution of content into a known frame rather than a fresh design. Nothing in
\S2 is investment advice, and no order is entered except against its checks.}
